\chapter{Introduction}
\label{ch:intro}

\section{Background and Motivation}
\label{sec:background}

Adaptive learning platforms remain dominated by static course graphs and pre-recorded lectures. While such material scales economically, it cannot react to the learner mid-sentence: it cannot recognise that a user mispronounced a word, that their gaze has drifted from the camera, or that a piece of code on the shared screen contains a bug. The recent generation of multimodal large language models (LLMs)---in particular Google's Gemini Live API, which streams audio, images and text in a single bidirectional channel---has changed this calculus. A single inference endpoint can now hold a spoken dialogue, watch the learner via webcam, look at the screen the learner shares, and reply with synthetic speech, all within latency budgets acceptable for tutoring ($\sim$300--700\,ms one-way).

Sprout is built on this premise. It is an AI-powered adaptive learning platform that personalises skill acquisition for adult learners, neurodiverse learners (with attention or autism-spectrum profiles), and vocational students whose practice naturally involves doing, demonstrating and being seen. The product surfaces three user-visible behaviours: natural-language skill discovery, visual knowledge maps (an AI generates a prerequisite graph of lessons in the React Flow canvas), and a multi-modal teaching assistant.

What makes those behaviours economically interesting is replacement pressure. Human one-on-one tutoring costs roughly \$20--\$50/hour; an LLM-driven session approaches \$2--\$3/hour at typical token usage. The technical question---how do we deliver such a session reliably, cheaply and observably?---is the subject of this report.

\section{Problem Statement}
\label{sec:problem}

The Sprout AI subsystem must satisfy four objectives concurrently:
\begin{enumerate}
    \item \textbf{Real-time multimodal communication.} A session must accept microphone audio, webcam video, screen frames and text within a single context window, and must stream the model's audio response back fast enough that turn-taking feels natural ($\leq$ 1\,s round-trip target).
    \item \textbf{Information recording.} Every dialogue turn, completion signal and feedback rating must be persisted to support context priming, parent and educator dashboards, and administrative audit, while respecting privacy constraints on raw media.
    \item \textbf{Direct exchange with the user.} The frontend must connect directly to the AI gateway---there must not be a Cloud Function in the audio path---so latency is dominated by the WAN leg only. Authentication must be enforced at the gateway, not by the model.
    \item \textbf{Operational predictability.} The system must support graceful failover, observable token consumption, an explicit maximum-session-duration policy, and rate limits to defend against abuse and to make freemium quotas tractable.
\end{enumerate}

A pre-existing Python/LiveKit pipeline that orchestrated GPT, Gemini and custom STT (Whisper) and TTS (Silero) plugins satisfied the first two objectives but failed the latter two: it held two server-side hops in the audio path, made authentication awkward, and produced opaque per-session cost figures. The present design replaces that pipeline with a thin gateway and is the focus of the report.

\section{Engineering Questions and Objectives}
\label{sec:questions}

The report is organised around three concrete engineering questions:
\begin{itemize}
    \item \textbf{Q1 (Latency \& topology).} Can a JWT-validated WebSocket proxy deliver lower median dialogue latency than the LiveKit-based agent system at equal or lower cost?
    \item \textbf{Q2 (Recording shape).} What Firestore schema and trigger graph minimise the number of round trips needed to (a) prime the AI tutor with previous context and (b) detect lesson completion across both agent-side and client-side observation?
    \item \textbf{Q3 (Pre-signup activation).} How can an anonymous user persist progress, accumulate value through full multimodal sessions, and later promote that progress into an authenticated account without losing data and without an account-merge race?
\end{itemize}

The corresponding engineering objectives are: (i) build a Cloud-Run-hosted proxy validating Firebase/Google JWTs and forwarding to Gemini Live; (ii) define an explicit Firestore data model; (iii) ship a typed TypeScript SDK; (iv) implement asynchronous knowledge-map generation via Cloud Tasks; (v) implement dual-channel lesson completion detection; and (vi) implement HTTP-only-cookie-based anonymous session linking.

\section{Contributions}
\label{sec:contributions}

The principal contributions of this report are:
\begin{enumerate}
    \item A reproducible system architecture from the browser WebSocket to the Gemini Live API and back, with explicit component boundaries and deployment artefacts.
    \item A formalisation of six core algorithms governing the AI subsystem.
    \item A Firestore-grounded data model that records every modality of dialogue plus aggregate session statistics.
    \item A typed client SDK contract (\texttt{@abyxolv/gemini-proxy-sdk}) that hides WebSocket reconnection, JWT refresh, and audio framing from the rest of the frontend.
    \item A forward path to multi-agent coordination based on Gemini function calling~\cite{midscene2024,sprout2026multiagent}.
\end{enumerate}

\section{Report Organisation}
\label{sec:organisation}

Section~\ref{ch:litreview} reviews related literature. Section~\ref{ch:context} describes the system context. Section~\ref{ch:methodology} formalises the six algorithms and the proposed multi-intent user representation. Section~\ref{ch:implementation} documents the implementation. Section~\ref{ch:evaluation} states the evaluation methodology. Section~\ref{ch:discussion} discusses achieved properties and limitations. Section~\ref{ch:operations} covers operations and reproducibility. Section~\ref{ch:conclusion} concludes with the multi-agent and function-calling roadmap.
